\section{Online partitions and the embedded comparator}
\label{sec:partitions}

Probabilistic tree embeddings use scale partitions to represent a metric
\cite{Bartal96,Bartal98,FRT04}. Local ball-mass ratios and their
logarithmic scale sums are familiar from measured descent
\cite[Lemma~1.7 and Section~1.3]{KLMN05}. Protection by heavy balls
and maintenance of separated nets also appear in Lee's evolving
fusible-HST construction \cite[Sections~3.1.2--3.2]{Lee18}; its
withdrawn competitive-ratio claim is discussed in
Section~\ref{sec:introduction}. Here the changing measure is the
posterior of the hidden comparator. The pilot potentials will pay
the edits induced by maintaining these partitions.

The partitions at distinct scales need not refine one another.  A point's
entire string of scale labels will specify its leaf in a fixed tree.  We
protect requests where the comparator measure is locally concentrated by
large disjoint balls.  At the other scales, chronological rosters of smaller
balls provide a separation estimate whose coefficient depends on the local
mass ratio.  The pilot potentials pay changes to both structures.

We use the finite metric, selected labeled comparator, scales, and filtration
of Section~\ref{sec:setup}.  In particular,
\[
 r_j=r_0\tau^{-j},\qquad r_0=\diam Y,\qquad
 40r_J<\min_{u\ne v}d(u,v),\qquad 1\le j\le J.
\]
All the numerical constants chosen below are absolute.  The constant
\(\tau\) also satisfies the absorption requirements of the allocation and
rounding arguments.  At a request \(x\), the true-metric posterior measures
obey
\begin{equation}\label{eq:partition-mover}
 \mu^+=\mu^--\nu+\delta_x,\qquad
 \nu(Y)=1,\qquad \nu\le\mu^-,\qquad
 e_x=\int d(x,y)\,d\nu(y).
\end{equation}
The selected mover is specified even on a zero move.  In particular every
ball centered at \(x\) has \(\mu^+\)-mass at least one, and
\(\E\sum_t e_{x_t}=Q\).

\begin{proposition}[Embedded comparator cost]\label{prop:embedded-cost}
There are finite alphabets \(\mathcal A_1,\ldots,\mathcal A_J\), depending
only on \(Y,k\) and the fixed constants, and online random maps
\[
 \xi_t:Y\longrightarrow\mathcal A_1\times\cdots\times\mathcal A_J
\]
with the following properties.  Each coordinate cell at level \(j\) has
diameter at most \(40r_j\) and an anchor within \(20r_j\) of all its points.
Thus \(\xi_t\) is injective.  Give the full alphabet tree the leaf distance
\(r_{j-1}\) when the first differing coordinate is \(j\).  For the selected
hidden members \(X_t^1,\ldots,X_t^k\),
\begin{equation}\label{eq:embedded-cost}
 P:=\E\sum_t\sum_{a=1}^k
 d_{\mathcal T}\bigl(\xi_{t-1}(X_{t-1}^a),\xi_t(X_t^a)\bigr)
 \le C\ell Q+D_{\rm part}.
\end{equation}
Here \(C\) is absolute and \(D_{\rm part}<\infty\) is independent of the
input law and its horizon.  A coordinate label common to successive maps
has a fixed anchor, except for the explicit heavy-ball reuse below, when
its anchor moves by at most \(40r_j\).
\end{proposition}

The maps for a given finite input law may use its true-prefix posteriors.
The alphabets and all initialization bounds will be fixed before that law
or its horizon is chosen.  Two forms of conditioning will be kept separate:
pilot potentials use filtering in \(\mathcal F_t\), whereas ball and
retirement probabilities average over the independent partition tape after
fixing the whole true input.

\subsection{Heavy balls and their edit cost}

Use the pilot template of Section~\ref{sec:pilot} with
\(\sigma=1\), \(R=512\), and Lipschitz bound one for \(g\) in scale units.
Let \(D_H=D_1=100R\).  Choose \(\gamma_H,\delta_H>0\) small enough for
Lemma~\ref{lem:pilot-local}, and also with
\(\gamma_H<1/100\), \(\delta_H<1/100\).  They may be decreased further
to meet the absolute requirement in Section~\ref{sec:travel}.
At scale \(r\), call the current request \emph{heavy} when
\begin{equation}\label{eq:heavy-condition}
 \mu^+(B(x,D_Hr))
 \le(1+\delta_H)\mu^+(B(x,\gamma_Hr)).
\end{equation}
Maintain an initially empty center set \(\mathcal H\).  If the request is
heavy and \(d(x,\mathcal H)>8r\), remove all centers within \(100r\) of
\(x\), and insert \(x\).  Otherwise leave \(\mathcal H\) unchanged.
Distance to the empty set is infinite.  Every inserted center receives a
fresh independent uniform radius in \([16r,20r]\), held until its removal.
These radii affect the balls and their labels, but never the center schedule.

Centers in \(\mathcal H\) have pairwise distance greater than \(100r\),
by induction on insertions.  Their balls are therefore disjoint.  If a new
ball intersects an old ball as subsets of \(Y\), reuse that old ball's
label.  Otherwise use a fresh label.  There is at most one intersected old
ball: two such old centers would be at distance at most \(80r\).  Moreover,
an intersected center is within \(40r\) of \(x\), so was removed in this
update.  The reused label's anchor consequently moves by at most \(40r\).
Intersection is meant in the finite metric itself; no geodesic or ambient
ball intersection is needed.  Heavy balls have precedence over every other
structure at the level.  In particular their final key regions are their
entire balls.

Write
\[
 E_x(r)=\{p\in Y:3r<d(p,x)\le120r\}.
\]
First we pay the mass in this annulus at heavy updates.  Let \(\rho_H\)
be zero on \([0,3]\), affine from zero to one on \([3,4]\), and one
thereafter, and set
\[
 g_{\mathcal H}(p)=\rho_H\bigl(d(p,\mathcal H)/r\bigr).
\]
At a heavy update, apply Lemma~\ref{lem:pilot-local} to \(\mu^+\) and
the old \(g\).  The inner mass is at least one and the shell ratio is at
most \(\delta_H\).  Every minimizing pilot configuration has a site
\(s\) with \(d(s,x)<r/4\) and \(z_s\ge1/2\), and no positive sites
at distances in \([r/4,256r]\) from \(x\).  The near site's old \(g(s)\)
is one, because every old center is more than \(8r\) from \(x\), and
its new value is zero.  For every \(p\in E_x(r)\),
\[
 2.75<d(s,p)/r\le120.25,
\]
so its surcharge bump \(a_s(p)\) equals one.  The decrease of this
surcharge is at least \(128r\mu^+(E_x(r))\).

An increase of \(g(p)\) requires removal of a center within \(4r\) of
\(p\), and cannot occur inside \(B(x,3r)\), where the new \(g\) is zero.
Thus every increase is supported at distances in \((3r,104r]\) from
\(x\).  All possible increasing-coefficient sites have zero weight by
the forbidden-site conclusion above.  Only the residual term can increase,
by at most \(r\mu^+(E_x(r))\), since its coefficient lies in \([0,1]\).
Using the old minimizer as a feasible candidate after the edit proves
\[
 \mathcal L_r(\mu^+,g_{\rm old})
 -\mathcal L_r(\mu^+,g_{\rm new})
 \ge127r\mu^+(E_x(r)).
\]
The all-step filtering and drift telescope of
Corollary~\ref{cor:pilot-edits} now gives
\begin{equation}\label{eq:heavy-annulus-budget}
 \E\sum_{t,j:\,\mathcal H_j\text{ updates at }t}
 r_j\mu_t^+(E_{x_t}(r_j))
 \le C\ell Q+D_H'.
\end{equation}
Here and below the constants denoted by \(D\) are finite initialization
bounds depending only on the fixed finite data.  Filtering and drift are
summed at the old parameters on every step, before testing whether an
update occurs; there is no conditioning of the filtering inequality on
the event of an update.

We will use two elementary scale sums repeatedly.  Set
\[
 L_j(x)=\log\frac{\mu^+(B(x,D_Lr_j))}
                         {\mu^+(B(x,\gamma_Lr_j))},
\]
where the absolute constants \(D_L>\gamma_L>0\) will be fixed below.
Lemma~\ref{lem:scale-logarithms}, with inner masses at least one and outer
masses at most \(k\), gives
\begin{equation}\label{eq:partition-log-sums}
 \sum_{j=1}^J L_j(x)\le p_L\log k,
 \qquad
 \#\{j:\text{the request is nonheavy at }j\}
 \le\frac{p_H\log k}{\log(1+\delta_H)}\le C\ell,
\end{equation}
where \(p_L,p_H\) are fixed positive integers satisfying
\(\tau^{p_L}\gamma_L\ge D_L\) and
\(\tau^{p_H}\gamma_H\ge D_H\).  For the second bound every nonheavy
level contributes more than \(\log(1+\delta_H)\) to its heavy-template
logarithmic ratio.  These sums involve one common measure \(\mu^+\) on
the step, so the scale telescope applies.

\begin{lemma}[Heavy edits]\label{lem:heavy-edits}
The expected sum of \(r_j\) times overall level-key changes caused by
heavy edits, evaluated at the hidden members' old physical positions with
the base structures held fixed, is at most \(C\ell Q+D_H''\).
\end{lemma}
\begin{proof}
Let \(\psi\) be the clipped increasing ramp on \([9,10]\).  For each
member at \(p\), add the auxiliary potential
\(r\psi(d(p,\mathcal H)/r)\).  Sum it over members and levels, and
telescope along actual hidden trajectories from one step end to the next.
Its initial value is at most \(k\sum_jr_j\) and it is nonnegative.
We bound an edit charge plus the full-step change of this potential.

Consider first a stationary member \(p\) at an update.  If
\(d(p,x)\le3r\) and an old center lies within \(10r\) of \(p\), both
the old and new heavy balls contain \(p\).  Their intersection forces
reuse, so its key does not change and its potential does not increase.
If no such old center exists, its old potential is \(r\) and its new
potential is zero, paying a possible charge of \(r\).  Outside this
inner ball, all affected old or new balls and all changing ramps lie in
\(B(x,120r)\): a removed center is within \(100r\), its ball has radius
at most \(20r\), and its ramp has range at most \(10r\).  The charge
plus potential increase is at most \(2r\), supported in \(E_x(r)\).

For the designated mover \(y\to x\), its final potential at an update
is zero.  The same argument pays its charge when \(d(y,x)\le3r\).
Otherwise the unpaid amount is at most \(r\), and
\(\sum_{r_j<d(x,y)/3}r_j\le C d(x,y)\).  At a step without an update
there is no edit charge.  On a heavy level a final center is within
\(8r\) of \(x\), so the final potential is again zero.  Only nonheavy
levels can have a positive physical-motion change; at each it is at most
\(d(x,y)\), by the ramp's Lipschitz bound.  Their number is bounded in
\eqref{eq:partition-log-sums}.

The conditional expected stationary counting measure equals
\(\mu^+-\delta_x\), and in particular is bounded by \(\mu^+\).
Consequently \eqref{eq:heavy-annulus-budget} pays the stationary errors,
and the other terms sum to \(C\ell Q\).  Nonnegativity and the stated
initial bound complete the telescope.
\end{proof}

\subsection{Chronological rosters and a fixed alphabet}

For the second pilot template choose \(\sigma>0\) so small that
\(3\sigma<1/8\) and \(10\sigma<\gamma_H\), and put \(R=4D_H\) and
\(D_L=100R\).  Choose \(0<\gamma_L<\sigma\) sufficiently small for
this template in Lemma~\ref{lem:pilot-local}.  The Lipschitz bound used
for its \(g\)'s will be eight in scale units.  In particular
\(\gamma_L<\gamma_H\) and \(D_L>D_H\).

Let
\[
 \mathcal H_k=\{1,2,4,\ldots,2^b\},\qquad
 b=\min\{n\in\N\cup\{0\}:2^n\ge\log k\}.
\]
The symbol \(\mathcal H_k\) denotes the set of tier parameters and is
distinct from the center set \(\mathcal H\).  Smaller tiers have
precedence.  A request qualifies in tier \(h\in\mathcal H_k\) at scale
\(r\) if
\begin{equation}\label{eq:tier-qualification}
 \mu^+(B(x,D_Lr))\le e^h\mu^+(B(x,\gamma_Lr)).
\end{equation}
The last tier always qualifies.  Fix an absolute \(C_K\), to be chosen
large below, and set \(K=K_h=\lceil e^{C_Kh}\rceil\).

Each tier keeps chronological insertion records and a random subroster
of balls used by the partition.  The insertion test below uses only
qualification and chronological ages.  Every entry young enough to
affect that test is deterministically present; later retirement changes
the partition but does not feed back into the insertion schedule.
This lets us estimate separation and retirement probabilities with that
schedule fixed.  The age cutoffs also bound the records relevant to the
test and partition.

Each pair \((j,h)\) has an initially empty chronological roster of
insertions.  The age of an entry centered at \(c\) is the number of
subsequent insertions in that same roster within distance \(20r\) of
\(c\).  Insert \(x\) if and only if it qualifies and there is no
chronological entry of current age less than \(K\) within \(r/2\) of
\(x\).  A noninsertion changes no ages.  These tests inspect the
chronological records, not their realized survival.  Records of age at
least \(K^2\) can be discarded forever: they cannot affect an insertion
test or any future presence.

The \emph{actual} roster used in the partition has the following lifetimes.
An entry is present deterministically at ages \(0,\ldots,K-1\).  At each
increment to an age in \(\{K,\ldots,K^2-1\}\), a still present entry
retires with fresh probability \(1/K\), independently.  It retires
compulsorily on reaching age \(K^2\).  At insertion its persistent radius
\(V\in[r,2r]\) has density
\begin{equation}\label{eq:tier-radius-density}
 f(v)=\frac{\lambda}{r}
       \frac{e^{-\lambda(v/r-1)}}{1-e^{-\lambda}}
       \one_{[r,2r]}(v),
 \qquad \lambda=\log(1+K^2).
\end{equation}
All entry radii and all lifetime trials are mutually independent.  Fresh
tape blocks are indexed by level, structure, and insertion ordinal, not
by the recyclable label slots.  The tier key at a point is the label of
its earliest present covering entry,
if any.  The overall level key is the heavy key if present, otherwise
the first applicable tier key, otherwise a fixed singleton key for that
point.  On a step, do the heavy edit first, followed by the roster edits
in increasing tier order; a roster's retirements and insertion form one
edit.  Each edit is evaluated at the old physical member positions.

Fixing the entire true input makes all comparator posteriors, heavy-center
schedules, roster insertion schedules, and chronological ages deterministic.
The schedules use neither radii nor lifetimes, nor allocation or rounding
states.  In particular conditioning on actual presences still leaves the
entry radii independent with law \eqref{eq:tier-radius-density}.

\begin{lemma}[Local age count]\label{lem:roster-count}
In a subset of \(Y\) of diameter at most \(20r\), there are at most
\(A\) chronological entries of age less than the positive integer \(A\).
In particular the bounds are \(K\) and \(K^2\) at the two age cutoffs,
regardless of how many entries have already retired.
\end{lemma}
\begin{proof}
If there were \(n>A\) such entries, the earliest of them would have
at least \(n-1\ge A\) later insertions within \(20r\) of its center,
contradicting its age.  Repeated centers cause no exception.
\end{proof}

When a request qualifies, its final tier covers \(x\) deterministically.
Indeed an insertion supplies an age-zero center at \(x\); if there is
no insertion, the test supplies an entry of age less than \(K\) within
\(r/2\), and no ages change.  That entry is present and its radius is
at least \(r\).  We shall refer to this as the tier's guaranteed coverage.

\begin{lemma}[Fixed labels and anchors]\label{lem:fixed-alphabet}
The construction can use fixed finite label pools at each level,
independent of the input law and horizon.  Labels can be assigned so that
distinct current records have distinct labels, including shadowed records,
and a freshly chosen label was absent throughout the preceding step and
all earlier intermediate states of the current step.  The sole exception
is the prescribed heavy-ball reuse.  Every point with a given current
label is within \(20r\) of its anchor; its cell has diameter at most
\(40r\).  A label common to consecutive maps has the same anchor or is
a reused heavy label with anchor displacement at most \(40r\).
\end{lemma}
\begin{proof}
Write \(n=|Y|\).  Applying Lemma~\ref{lem:roster-count} to each singleton
shows that a tier has at most \(M=nK^2\) chronological records of age
less than \(K^2\).  Reserve a record's label until that age cutoff,
even if its actual lifetime has ended or its ball is completely shadowed
by earlier structures.  This convention bounds all reserved labels by
\(M\); it does not alter the partition.  During a step at most one new
record is inserted.  The records present at any intermediate time during
two consecutive steps belong to their two starting record sets and at
most two new records, so their total number is at most \(2M+2\).
A pool of \(4M+4\) slots suffices to choose a fresh slot outside this
entire set whenever needed.  A slot released on the previous step is
still excluded on the current step.  This is a finite implementation of
the stated cooldown rule, not a growing list of all historical labels.

The heavy structure has at most \(n\) centers and one insertion per step,
so a separate pool of \(4n+4\) slots works by the same rule, with reuse
allowed precisely for the unique intersected ball.  Use disjoint pools
for different tiers and the heavy structure, and \(n\) additional fixed
singleton labels.  Pool sizes are chosen for every level in advance and
do not depend on the realized center schedule.  A chronological record
that persists keeps its label and center even while shadowed.  Thus any
common slot can only change its anchor by the explicit heavy reuse.

Heavy and tier keys are anchored at their centers, and singleton keys at
their points.  Every key cell is a subset of its defining ball, whose
radius is at most \(20r\); restricting by precedence only shrinks it.
The diameter and anchor bounds follow.  Heavy precedence further implies
that a heavy key's region is its entire ball.  Empty key regions cause
no difficulty; their reserved anchors and slots are still defined.
\end{proof}

\subsection{Separation in the final map}

\begin{lemma}[First hit]\label{lem:first-hit}
Fix a true input, a roster at a fixed step, and points \(x,y\) with
\(d=d(x,y)<r/4\).  The probability that the earliest present ball
covering either point covers exactly one is at most \(Chd/r\).
The expectation is over the partition tape.
\end{lemma}
\begin{proof}
Condition on all presences, leaving the radii independent.  For a center
\(c\), put \(a=\min\{d(c,x),d(c,y)\}\) and
\(p_a=\Pr(V\ge a)\).  The two distances differ by at most \(d\), so
separation requires \(V\in[a,a+d]\), up to endpoints of probability
zero.  If \(r\le a\le2r\), direct integration of
\eqref{eq:tier-radius-density} gives
\[
 f(a)=\frac{\lambda}{r}
       \left(p_a+\frac1{e^\lambda-1}\right).
\]
The density is decreasing, hence
\begin{equation}\label{eq:radius-separation-tail}
 \Pr\bigl(V\in[a,a+d]\bigr)
 \le\frac{\lambda d}{r}
       \left(p_a+\frac1{e^\lambda-1}\right).
\end{equation}
For \(a<r\), use the maximum density at \(r\) and \(p_a=1\);
for \(a>2r\), the left side is zero.  Thus the same bound holds when
an interval crosses a support endpoint.

List the present balls in chronological order and let \(q_i\) be the
probability all balls preceding ball \(i\) miss both points.  The event
that ball \(i\) hits at least one has probability \(p_{a_i}\), so
independence gives
\[
 \sum_i q_i p_{a_i}
 =\Pr(\text{some present ball hits }x\text{ or }y)\le1.
\]
Every potentially hitting center lies within \(2r\) of one of the
points.  Those centers have diameter at most \(4r+d<5r\), and all
present entries have age less than \(K^2\).  There are therefore at
most \(K^2\) candidates by Lemma~\ref{lem:roster-count}.  Multiplying
\eqref{eq:radius-separation-tail} by \(q_i\) and summing gives
\[
 \Pr(\text{separating first hit})
 \le\frac{\lambda d}{r}
       \left(1+\frac{K^2}{e^\lambda-1}\right)
 =\frac{2\lambda d}{r}.
\]
Finally \(\lambda=\log(1+K^2)\le Ch\), since \(C_K\) is an absolute
constant and \(h\ge1\).  The bound is uniform in the conditioned
presences, so it also holds without that conditioning.
\end{proof}

Let \(\xi_{t,j}\) denote the final overall key at level \(j\).
For every fixed input and every \(y\in Y\),
\begin{equation}\label{eq:final-map-separation}
 \E_{\rm tape}\sum_{j=1}^Jr_j
 \one_{\{\xi_{t,j}(x)\ne\xi_{t,j}(y)\}}
 \le C\ell d(x,y).
\end{equation}
To prove this, put \(d=d(x,y)\); the case \(d=0\) is immediate.
The radii \(r_j\le4d\) contribute at most \(C d\) by geometric
summation.  At a larger heavy level a final center lies within \(8r\)
of \(x\), so its ball of radius at least \(16r\) contains both points.
There is agreement regardless of its sampled radius.

At a larger nonheavy level, the centers of heavy balls that could contain
one of the two points lie within \(20r\) of their union, a set of
diameter at most \(40r+d<41r\).  There is at most one such center.
Its uniform radius separates the two points with probability at most
\(d/(4r)\).  Next let \(h_*(j)\) be the first qualifying tier at this
level.  It supplies guaranteed coverage of \(x\).  If no heavy ball
separates the points and no tier through \(h_*(j)\) has a separating
first hit, their overall keys agree: the first structure hitting either
hits both, and such a structure exists by guaranteed coverage.
Lemma~\ref{lem:first-hit} and a union bound therefore give probability
at most
\[
 \frac{Cd}{r}\left(1+\sum_{h\le h_*(j)}h\right)
 \le\frac{Cd}{r}(1+L_j(x)).
\]
Indeed the dyadic sum is less than \(2h_*(j)\), and the first dyadic
value at least \(L_j(x)\) is at most \(2\max\{1,L_j(x)\}\).
After multiplying by \(r\), sum the additive ones only over nonheavy
levels and the logarithms by \eqref{eq:partition-log-sums}.  This proves
\eqref{eq:final-map-separation}.  In particular the argument does not
multiply the number of nonheavy levels by a worst-case tier coefficient;
the individual coefficients telescope through the local mass ratios.

\subsection{Noncoverage and the cost of roster edits}

\begin{lemma}[Tier edits]\label{lem:tier-edits}
The expected sum of \(r_j\) times overall level-key changes due to
roster edits, evaluated at the old physical positions in the stated edit
order, is at most \(C\ell Q+D_L'\).
\end{lemma}
\begin{proof}
For each level and tier use the nonnegative potential
\[
 \Phi_{j,h}=r_j\sum_{a=1}^k
 \one_{\{X^a\text{ is not covered by that tier}\}},
\]
ignoring all other structures in the definition of coverage.  Telescope
from step end to step end, so the potential includes physical movement.
The sum of initial values is at most
\(k|\mathcal H_k|\sum_jr_j\), an allowed fixed initialization bound.

If there is no insertion, the roster has no edit.  Stationary members
give zero potential change, and only the mover \(y\to x\) can increase
the potential.  This requires \(y\) covered and \(x\) uncovered.  It
is impossible in a qualified tier, by guaranteed coverage.  In an
unqualified tier, Lemma~\ref{lem:first-hit} gives expected positive
increase at most \(Chd(x,y)\) when \(r>4d(x,y)\).  For
\(r\le4d(x,y)\), the bound follows from the trivial increase at most
\(r\le4h d(x,y)\).  Unqualification says \(h<L_j(x)\), and
\begin{equation}\label{eq:unqualified-tier-sum}
 \sum_{j=1}^J\sum_{h\in\mathcal H_k:\,h<L_j(x)}h
 \le2\sum_{j=1}^J L_j(x)\le C\ell.
\end{equation}
Thus all noninsertion increases cost at most \(C\ell Q\).

At an insertion, let \(C_a\in\{0,r\}\) be a member's overall key
charge caused by this tier edit at its old position.  Let \(U_-\) be
old tier noncoverage at that position and \(U_+\) final tier
noncoverage at its final physical position.  In particular \(U_+=0\)
for the mover, because the newly inserted ball covers \(x\).  The
following table accounts for the entire quantity
\(C_a+r(U_+-U_-)\).
\begin{center}
\small
\begin{tabular}{p{0.31\linewidth}p{0.30\linewidth}p{0.28\linewidth}}
\toprule
Old tier status & Stationary member & Designated mover \\
\midrule
First covering entry persists &
\(C_a=0\), \(U_-=U_+=0\) &
\(C_a=0\), \(U_-=U_+=0\) \\
\addlinespace
No covering entry &
\(C_a\le r(1-U_+)\), \(U_-=1\); total at most zero &
\(C_a\le r\), \(U_-=1\), \(U_+=0\); total at most zero \\
\addlinespace
First covering entry retires &
\(C_a\le r\), \(U_-=0\), \(U_+\le1\); total at most \(2r\) &
\(C_a\le r\), \(U_-=U_+=0\); total at most \(r\) \\
\bottomrule
\end{tabular}
\end{center}
For the first row, new entries are appended after all existing ones, so
a surviving first coverer keeps the tier key unchanged.  Other structures
are held fixed during this edit.  For the second row, an overall key can
change only if this previously uncovered point becomes covered by the
tier.  These facts also hold when an earlier structure shadows the tier,
in which case its charge may simply be zero.

There is a further zero-error case at a heavy level.  If the old position
lies in \(B(x,\gamma_Hr)\), the final heavy ball already covers it
during all tier edits.  Hence \(C_a=0\).  The new tier ball covers
that position if it is stationary, and covers the final point \(x\)
if it is the mover.  Thus \(U_+=0\) and no retirement error is needed
in this inner ball.

We next bound the probability of the exceptional retirement row.  Fix
the full true input and a point \(p\).  Conditional on the old realized
presences and radii, a pre-edit first coverer at a noncompulsory age can
be removed on this step only by its fresh independent trial, of
probability at most \(1/K\).  For compulsory removal use an
unconditional bound over the tape.  At most \(K^2\) chronological
candidates of age less than \(K^2\) have centers within \(2r\) of
\(p\).  A candidate reaching age \(K^2\) must previously have survived
the \(K^2-K\) trials at ages \(K,\ldots,K^2-1\).  Consequently
\begin{equation}\label{eq:actual-retirement-bound}
 \Pr_{\rm tape}(\text{the old first coverer of }p\text{ retires})
 \le\frac1K+K^2(1-1/K)^{K^2-K}\le\frac{C}{K}.
\end{equation}
The last inequality follows, for example, from
\(K^2(1-1/K)^{K^2-K}\le K^2e^{-(K-1)}\le C/K\).
If \(d(p,x)>22r\), no covering entry can have its age incremented:
its center is within \(2r\) of \(p\) but would have to be within
\(20r\) of \(x\).  The retirement probability is then zero.
The compulsory part of \eqref{eq:actual-retirement-bound} is not asserted
conditional on an observed surviving old roster; it averages over the
independent tape with the chronological schedule fixed.

For an insertion at time \(t\), level \(j\), and tier \(h\), write
\begin{equation}\label{eq:tier-insertion-mass}
 W_{t,j,h}=
 \begin{cases}
 \mu^+\{p:\gamma_Hr<d(p,x)\le D_Hr\},&\text{if heavy},\\
 \mu^+(B(x,D_Lr)),&\text{if nonheavy}.
 \end{cases}
\end{equation}
Although the value at a fixed step and level does not depend on \(h\),
the indices specify the insertion being charged.  The stationary errors
are confined to this region, since \(D_H>22\), and the heavy inner ball
has zero error.  Average \eqref{eq:actual-retirement-bound} first over
the tape with the full input fixed, and then over input continuations
given the true prefix.  The insertion event and \(W_{t,j,h}\) are
fixed by that prefix.  The expected stationary counting measure is
\(\mu^+-\delta_x\le\mu^+\).  Thus the expected stationary error is
at most \(CrW_{t,j,h}/K\).  This averaging statement is conditional
only on the true prefix, not on the complete partition tape in
\(\mathcal F_t\).

At a nonheavy insertion, the mover's error is also covered by
\(CrW_{t,j,h}/K\), since \(W_{t,j,h}\ge\mu^+(\{x\})\ge1\).
At a heavy insertion, its remaining error is at most \(Cr/K\) and
only occurs if \(d(y,x)>\gamma_Hr\).  For each tier,
\[
 \sum_{r_j<d(x,y)/\gamma_H}\frac{Cr_j}{K_h}
 \le\frac{C d(x,y)}{K_h},\qquad
 \sum_{h\in\mathcal H_k}\frac1{K_h}\le C.
\]
Here the fixed factor \(1/\gamma_H\) is absorbed into \(C\).
These mover errors therefore cost at most \(CQ\).  It remains to prove
\begin{equation}\label{eq:insertion-charge-budget}
 \E\sum_{t,j,h:\,\text{insertion}}\frac{r_jW_{t,j,h}}{K_h}
 \le C\ell Q+D_{\mathrm{ins}}.
\end{equation}

\paragraph{An auxiliary short roster.}
We need a potential drop at an insertion to pay the remaining insertion charge.
The insertion test excludes young records near $x$, but an older surviving
ball may still cover $x$.  A potential built from the actual surviving
centers would therefore have no guaranteed change there.  Instead we use
an auxiliary roster containing only entries of age less than $K$.
The insertion test then forces a vacant neighborhood of $x$ before the
new entry arrives.  Random retirement before the age cutoff keeps the
expected increases elsewhere small.

This auxiliary roster uses the same deterministic insertion records and
ages.  An entry starts present at
age zero, has independent retirement probability \(1/\sqrt K\) at
each subsequent age increment, and retires compulsorily at age \(K\).
Its sole output is the deterministic expected function below; it does
not enter the partition.  Let \(\rho_L\) be the clipped increasing ramp on
\([1/8,1/4]\), taking value one at infinity, and set
\begin{equation}\label{eq:short-roster-g}
 g(p)=\E_{\rm short}
       \rho_L\bigl(d(p,\mathcal S)/r\bigr),
\end{equation}
where \(\mathcal S\) is the set of surviving hypothetical centers.
The expectation is over these auxiliary lifetimes, given the
chronological schedule.  The function \(g\) lies in \([0,1]\) and
has \(r\Lip(g)\le8\); it is a deterministic function of the true
prefix, and is unchanged at noninsertions.

At an insertion the test excludes all old centers of age less than
\(K\) within \(r/2\) of \(x\).  Every hypothetical survivor has age
less than \(K\).  Hence before insertion \(g=1\) throughout
\(B(x,r/8)\), while the new age-zero entry gives \(g=0\) there afterward.
Elsewhere couple old and new hypothetical lifetimes by continuing the
old trials.  If the old ramp at \(p\) can increase, a nearest surviving
center within \(r/4\) of \(p\), selected by a fixed tie rule, must
retire.  Its noncompulsory fresh trial has probability \(1/\sqrt K\).
There are at most \(K\) possible compulsory candidates within \(r/4\)
of \(p\) by Lemma~\ref{lem:roster-count}, and a candidate reaching age
\(K\) must have survived \(K-1\) previous trials.  Thus
\begin{equation}\label{eq:short-retirement-bound}
 (g_{\mathrm{new}}(p)-g_{\mathrm{old}}(p))_+
 \le\varepsilon_K:=K^{-1/2}+K(1-K^{-1/2})^{K-1}
 \le C K^{-1/2}.
\end{equation}
The last bound follows by comparing the second term with
\(K\exp(-(K-1)/\sqrt K)\); its ratio to \(K^{-1/2}\) is bounded
over integers \(K\ge2\).  Only centers within \(20r\) of \(x\)
change age.  Therefore positive changes of \(g\) are supported in
\begin{equation}\label{eq:short-positive-support}
 B(x,21r)\setminus B(x,r/8).
\end{equation}
The auxiliary estimate is $O(K^{-1/2})$, whereas the bound used for
actual retirements is $O(K^{-1})$.  The exponentially large choice
of $K_h$ will make the auxiliary error small enough compared with
the insertion gain.

\paragraph{The insertion gain in the pilot.}
Fix an insertion, abbreviate \(W=W_{t,j,h}\), and let \(z\) be an old
minimizer of the second-template pilot at \(\mu^+\) and
\(g_{\mathrm{old}}\).  Its residual
\(R_z(p)=1-\sum_s z_s b_s(p)\) lies in \([0,1]\).  The change of
its objective when \(g\) is edited is exactly
\begin{equation}\label{eq:tier-candidate-change}
 r\int\left[
   \Delta g(p)R_z(p)+256\sum_s z_s\Delta g(s)a_s(p)
 \right]d\mu^+(p).
\end{equation}
At a nonheavy level, its positive part is at most
\(CrW/\sqrt K\).  Indeed the residual increase is supported in
\(B(x,21r)\).  Any increasing-coefficient site lies in that same
ball, and its surcharge support lies in
\(B(x,(21+2R)r)\subset B(x,D_Lr)\).  At any test point the weights
of all sites with nonzero \(a_s\) sum to at most one, by the pilot
capacity constraint and \(2R<D_*\).  Now apply
\eqref{eq:short-retirement-bound} to both terms of
\eqref{eq:tier-candidate-change}.

At a heavy level the same bound holds with the annular definition of
\(W\), rather than the possibly much larger inner mass.  To see this,
every possible increasing-coefficient site satisfies
\(r/8<d(s,x)\le21r\).  Since
\(1/8-\gamma_H>\sigma\) and \(21+\gamma_H<R\), its bump \(a_s\)
equals one throughout \(B(x,\gamma_Hr)\).  Its \(b_s\)-support lies
inside \(B(x,D_Hr)\).  Writing
\(M_H=\mu^+(B(x,\gamma_Hr))\), the coefficient of \(z_s\) in the
old objective, divided by \(r\), is at least
\[
 256M_H-2\mu^+(B(x,D_Hr))
 \ge[256-2(1+\delta_H)]M_H>0.
\]
Deleting that weight preserves both capacity and separation constraints,
so optimality forces \(z_s=0\).  Hence no surcharge increases in
\eqref{eq:tier-candidate-change}; the residual increases only in
\eqref{eq:short-positive-support}, a subset of the heavy annulus in
\eqref{eq:tier-insertion-mass}.  Its positive contribution is again
at most \(CrW/\sqrt K\).

It remains to exhibit the negative contribution.  Put
\(M_L=\mu^+(B(x,\gamma_Lr))\).  Qualification and the definition
of \(W\) give, in both the heavy and nonheavy cases,
\begin{equation}\label{eq:insertion-inner-mass}
 M_L\ge e^{-h}\mu^+(B(x,D_Lr))\ge e^{-h}W.
\end{equation}
There are two exhaustive cases.

If some site \(s\) within \(3\sigma r\) of \(x\) has
\(z_s\ge1/2\), then \(g(s)\) drops from one to zero because
\(3\sigma<1/8\).  On the entire heavy annulus
\(\gamma_Hr<d(p,x)\le D_Hr\), the bump \(a_s(p)\) equals one:
the lower distance is greater than \(\gamma_H-3\sigma>\sigma\)
in scale units, and the upper distance is less than
\(D_H+3\sigma<R\).  Its surcharge drop is at least \(128rW\)
when the request is heavy.  When it is nonheavy, failure of
\eqref{eq:heavy-condition} says that this annulus has mass greater than
\(\delta_H M_H\ge\delta_H M_L\).  The drop is then at least
\(128\delta_HrM_L\ge128\delta_Hre^{-h}W\).

If no such site exists, then for every
\(p\in B(x,\gamma_Lr)\), any contributing \(b_s(p)\) has its
center within \((2\sigma+\gamma_L)r<3\sigma r\) of \(x\), and
has weight less than \(1/2\).  At most one positive site contributes
there, since two centers within \(2\sigma r\) of the same point
would violate their \(5\sigma r\) separation.  It follows that
\(R_z(p)\ge1/2\) on this inner ball.  The unit drop of \(g(p)\)
there gives a residual decrease at least
\(rM_L/2\ge re^{-h}W/2\).

Combining either negative contribution with the established positive
error, and using the old optimum as a candidate for the new minimum,
gives absolute \(c,C>0\) such that
\begin{equation}\label{eq:tier-pilot-drop}
 \mathcal L_r(\mu^+,g_{\mathrm{old}})
 -\mathcal L_r(\mu^+,g_{\mathrm{new}})
 \ge rW\left(ce^{-h}-\frac C{\sqrt K}\right).
\end{equation}
Choose \(C_K\) sufficiently large, once and for all, that for every
\(h\ge1\),
\[
 C e^{-C_Kh/2}\le(c/2)e^{-h},\qquad
 e^{-C_Kh}\le(c/2)e^{-2h}.
\]
Since \(K\ge e^{C_Kh}\), multiplying
\eqref{eq:tier-pilot-drop} by \(e^{-h}\) yields
\begin{equation}\label{eq:weighted-insertion-drop}
 e^{-h}\bigl[
 \mathcal L_r(\mu^+,g_{\mathrm{old}})
 -\mathcal L_r(\mu^+,g_{\mathrm{new}})\bigr]
 \ge(c/2)e^{-2h}rW\ge\frac{rW}{K}.
\end{equation}

For each tier apply Corollary~\ref{cor:pilot-edits} with its held
function \eqref{eq:short-roster-g} and fixed weight \(e^{-h}\).
The template and Lipschitz constants are the same for all tiers, and
\(\sum_{h\in\mathcal H_k}e^{-h}\le\sum_{n\ge1}e^{-n}<\infty\).
Thus the filtering and summed drift budgets remain \(C\ell Q\),
with no factor equal to the number of tiers.  At noninsertions there
is no parameter edit.  At insertions
\eqref{eq:weighted-insertion-drop} pays exactly the desired charges.
The potentials are nonnegative and their total initial value is at most
\(2k\sum_jr_j\sum_he^{-h}\).  This proves
\eqref{eq:insertion-charge-budget}.

Finally sum the per-member case table over all insertions, add the
noninsertion estimate \eqref{eq:unqualified-tier-sum}, and telescope the
noncoverage potentials.  The expected retirement errors are covered by
\eqref{eq:insertion-charge-budget} and the heavy-mover geometric sum.
Their terminal nonnegative values can be discarded.  This proves the
lemma with an initialization bound independent of the law and horizon.
\end{proof}

\subsection{The fixed tree and its posterior interface}

\begin{proof}[Proof of Proposition~\ref{prop:embedded-cost}]
Let \(\mathcal A_j\) be the fixed level alphabet in
Lemma~\ref{lem:fixed-alphabet}, and take the tree of all prefixes of
words in \(\mathcal A_1\times\cdots\times\mathcal A_J\), including
unused words.  Initially the center structures are empty and the map
uses singleton labels.  At every step map a point to its complete
string of final keys.  The last-level cell diameter is at most
\(40r_J\), less than the minimum nonzero distance of \(Y\); hence
last-level cells are singletons, giving injectivity and a requested
leaf with anchor exactly equal to the actual request.

For any two words \(u,v\),
\begin{equation}\label{eq:key-to-tree-charge}
 d_{\mathcal T}(u,v)
 \le\tau\sum_{j=1}^Jr_j\one_{\{u_j\ne v_j\}}.
\end{equation}
Indeed, if their first disagreement is \(j\), the left side is
\(r_{j-1}=\tau r_j\).  For each member on each step, first edit all
level keys at its old physical position, and then move it to its new
position under the final map.  The triangle inequality for key
disagreements bounds the right side of
\eqref{eq:key-to-tree-charge} by the heavy edit charges, all tier edit
charges, and the final-map physical disagreements.  The first two have
budgets \(C\ell Q+D\) by Lemmas~\ref{lem:heavy-edits} and
\ref{lem:tier-edits}.  Only the designated mover has a nonzero final
physical disagreement; \eqref{eq:final-map-separation}, applied after
fixing the true input and hence that member's old position, bounds its
expected charge by \(C\ell d(x,y)\).  Averaging over inputs and
summing proves \eqref{eq:embedded-cost}; \(\tau\) is absolute.

All initialized potentials above are bounded by fixed multiples of
\(k(1+|\mathcal H_k|)\sum_jr_j\).  Their number of time updates never
enters the bound.  The finite pools and the full prefix tree likewise
depend only on \(Y,k\) and the chosen absolute constants.  Consequently
the additive constant is uniform over input laws and horizons.
\end{proof}

For completeness, the coupling used by the allocation theorem is now
fully specified.  Put
\[
 N_v(t)=\#\{a:\xi_t(X_t^a)\text{ has prefix }v\},\qquad
 m_v(t)=\E[N_v(t)\mid\mathcal F_t].
\]
The final map and the requested leaf are available from the observed true
prefix and independent partition tape before virtual service.  The
requested leaf contains the designated comparator member after drift,
with conditional probability one.  The filtering substep at time \(t\)
uses \(N_v(t-1)\) with the \emph{old} map \(\xi_{t-1}\), comparing
its laws given \(\mathcal F_{t-1}\) and \(\mathcal F_t\).  The drift
substep compares that same old-map variable with \(N_v(t)\), both
conditioned on \(\mathcal F_t\), along the hidden labeled coupling
whose expected cost is \eqref{eq:embedded-cost}.  In particular the
allocation drift budgets are paid by this \(P\).  No ball-separation
or survival probability is asserted after conditioning on the complete
tape; those estimates have already been averaged to establish \(P\).

For a finite horizon the partition outcomes can, if desired, be encoded
on a finite probability space: on finite \(Y\), a radius is used only
through its ball subset, so its continuous law has finitely many relevant
outcomes, and only finitely many insertions and lifetime trials occur.
The expected short-roster functions require no additional random
observations.  The construction therefore meets the finite-data
requirements of Theorem~\ref{thm:allocation}.

Apply that theorem once to the fixed full tree.  Its trackers continue
at every vertex when a label becomes unused and later reappears; they
are not restarted when a slot is recycled.  The resulting bound is
\(V\le C\ell P+D_1\) with a single fixed-tree initialization bound.
Partition schedules use only the true-prefix posteriors and
chronological ages, so neither these allocations nor the independent
rounding randomness feed back into the partition construction.
